\begin{figure}[t]
\centering
\figplaceholder[1.4in]{\columnwidth}{Fig.~5 --- ablation effects by task}{%
\textbf{Layout:} column width, two heatmaps side by side with shared row labels: left \emph{native latency}, right \emph{+150\,ms}.
\textbf{Rows (top to bottom):} no target history; no future supervision; no scene-flow supervision; no dedicated SE(3); uniform-fidelity forecast.
\textbf{Columns:} All; the five families (Grasp, Place, Pass, Capture, Impact); a thin separator; then two task groups fixed in advance from Table~\appref{tab:roles}: \emph{rotation-sensitive} (tasks whose terms include $r\theta$) and \emph{context-active} (tasks with an obstacle clearance $m$, such as the rotating wire, or a driver: belt, ramp, striker).
\textbf{Cells:} $\Delta$SR = control minus full model in points, mean of three training seeds, printed with one decimal; diverging color centered at 0 (\textcolor{failcol}{red} = drop, blue = gain), one symmetric scale for both heatmaps, saturated beyond $\pm$15 points; bold value when the paired 95\% bootstrap interval excludes 0.
\textbf{Rules:} same held-out scenes for all rows; group membership listed in the appendix; no t3\_phase exclusion here (sensitivity analysis in the appendix).}
\caption{\textbf{Ablation effects by task.} SR difference of each control to the full
model (points) with native latency (left) and 150\,ms of added latency (right), by
task family and by two task groups fixed from Table~\appref{tab:roles}: tasks whose
error depends on rotation ($r\theta$) and tasks whose context constrains or drives
the target. Bold: the paired 95\% interval excludes zero. \fillin{Summarize where
each component matters.}}
\label{fig:analysis}
\end{figure}